% ---
% id: meat-core/thm:removal
% kind: thm
% title: "Sound removal of a state-blind sector"
% summary: >-
%   A state-blind sector may be deleted when a mixture of states on distinct
%   other sectors, all key-inert, reproduces its responses and its constraint
%   values. Deleting it changes neither the value of the program nor its
%   attainment.
% prominence: standard
% status: implemented
% uses: [def:feasible, def:maps, prop:exact-cancellation, cor:f-nonpos, def:classes]
% assumes: []
% refs: []
% provenance: newly-derived
% implements: certiqs_key.meat_core.declaration.sectors:verify_removal_witness
% lean:
% notes: "A party-local value is an operator on the other measured factors' blocks, which the move does not preserve in general, so the theorem excludes it."
% ---

A state-blind sector answers every state alike, so its weight can be moved
onto other sectors that reproduce its answers. When those sectors carry no
key entropy, the move removes one non-positive deficit and changes nothing
else.

\begin{theorem}\label{thm:removal}
Let $\CList$ carry no entry declared \tref{party-local}, let
$s^{*}\in\SecSet$ be state-blind, and suppose there exist
\emph{pairwise distinct} sectors
$s_1,\dots,s_L\in\SecSet\setminus\{s^{*}\}$, states
$\chi_i\in\Seq{\Qsys_{s_i}}$, and weights $\mu_i\ge0$ with
$\sum_{i=1}^{L}\mu_i=1$, such that
\begin{enumerate}[label=(RM\arabic*),leftmargin=3.0em]
\item \emph{response match:}\quad
$c_{v,s^{*}}=\sum_{i}\mu_i\tr\big[F^{(s_i)}_v\chi_i\big]$ for every
$v\in\Vset$;
\item \emph{key-inert targets:}\quad every pair $(\hat c,s_i)$ with
$\hat c\in\Chkey$ is key-inert;
\item \emph{constraint match} (vacuous for $\CList=\varnothing$):\quad
$o_{k,s^{*}}=\sum_i\mu_i\tr\big[\hat O^{(s_i)}_k\chi_i\big]$ for every
$k$.
\end{enumerate}
Then restricting the infimum \eqref{eq:hlow-def} to
$\{\omega\in\Feas:\omega_{s^{*}}=0\}$ changes neither its value nor its
attainment: the sector $s^{*}$ may be deleted from the variable list, and
its rows deleted from the generated maps, without loss.
\end{theorem}

\begin{proof}
Given $\omega\in\Feas$ and $\mathbf q\in\acSet$, let
$m\coloneqq\tr_{\Qsys_{s^{*}}}\omega_{s^{*}}\in\Pos{A}$ and define
$\hat\omega$ by
\begin{equation*}
  \hat\omega_{s^{*}}\coloneqq0,\qquad
  \hat\omega_{s_i}\coloneqq\omega_{s_i}+\mu_i\,m\otimes\chi_i,\qquad
  \hat\omega_{s}\coloneqq\omega_{s}\ \ (s\notin\{s^{*},s_1,\dots,s_L\}).
\end{equation*}
The three assignments are mutually consistent, which the two index
conditions of the hypothesis ensure: no $s_i$ equals $s^{*}$, so
the first two do not collide, and the $s_i$ are pairwise distinct, so the
second assigns one value to each block it touches.
\emph{Feasibility.} Each block is PSD; sector-block-diagonality is
manifest. With Alice sectors, $\omega_s=\bigoplus_t\omega^{(t,s)}$ is block
diagonal along $\mathbb{T}$, hence so is $m=\bigoplus_tm_t$ with
$m_t=\tr_{\Qsys_{s^{*}}}\omega^{(t,s^{*})}$, and so is every
$m\otimes\chi_i$: the modified point is block diagonal along the
\tref[joint sector]{joint sectors}, and the move takes place inside each
Alice sector separately. \emph{Normalization} --- the condition that holds
in both source modes, and the only one in
\tref{untrusted-source mode}:
$\tr\hat\omega=\tr\omega-\tr m+\sum_i\mu_i\,\tr m\,\tr\chi_i
=\tr\omega=1$, using $\tr\chi_i=1$, $\sum_i\mu_i=1$. \emph{Alice's
marginal}, in \tref{trusted-source mode} only: the same computation at the
level of the partial trace,
$\tr_{\Qsys}\hat\omega
=\tr_{\Qsys}\omega-m+\sum_i\mu_i\,m\,\tr\chi_i
=\tr_{\Qsys}\omega=\tau_A$; the normalization statement is its trace,
so replacing the marginal by normalization is exactly the step that
carries this theorem into untrusted-source mode.
\emph{Constraints.} The $s^{*}$ contribution to
$\tr_{\Qsys}[(\idon{A}\otimes\hat O_k)\omega]$ is
$\tr_{\Qsys_{s^{*}}}[(\idon{A}\otimes o_{k,s^{*}}\idon{})\,\omega_{s^{*}}]
=o_{k,s^{*}}m$, and the added contribution is
$\sum_i\mu_i\tr[\hat O^{(s_i)}_k\chi_i]\,m=o_{k,s^{*}}m$ by (RM3); all
other sectors unchanged. In untrusted-source mode the constraint value
is the trace of that same operator, so the identity gives
$o_{k,s^{*}}\tr m$ on both sides.
\emph{Statistics.} For every $(x,v)$, the $s^{*}$ term of
\eqref{eq:Phimap} is
$\tr[(M_x\otimes c_{v,s^{*}}\idon{})\,\omega_{s^{*}}]
=c_{v,s^{*}}\tr[M_xm]$, and the added terms give
$\sum_i\mu_i\tr[M_xm]\tr[F^{(s_i)}_v\chi_i]=c_{v,s^{*}}\tr[M_xm]$ by
(RM1). Hence $\Phimap[\hat\omega]=\Phimap[\omega]$, resolved per outcome
pair, so every $P_{\hat c}$ is unchanged.
\emph{Objective.} The divergence terms depend on
$(\mathbf q,\Phimap[\cdot])$ only, hence are unchanged. For
$K_{\hat c}$: by (RM2) and Proposition~\ref{prop:exact-cancellation},
$f_{\hat c,s_i}\equiv0$ on \emph{every} input, before and after the
modification; $f_{\hat c,s}$ for other retained $s$ is unchanged; and the
term $f_{\hat c,s^{*}}(\omega_{s^{*}})\le0$
(Corollary~\ref{cor:f-nonpos}) is deleted. Hence
$K_{\hat c}(\hat\omega)\ge K_{\hat c}(\omega)$, and by monotonicity of the
kernel $d$ in its second argument,
$O_{\alpha}(\mathbf q,\hat\omega)\le O_{\alpha}(\mathbf q,\omega)$.
Therefore the restricted infimum is no larger than the unrestricted one;
being a restriction it is also no smaller; and the map
$\omega\mapsto\hat\omega$ carries minimizers to restricted minimizers.
\end{proof}
